
RAG systems are often described as a simple two-stage pipeline: retrieve documents, then generate an answer. Robustness failures show that this description is too coarse. In practice, a deployed RAG system includes corpus acquisition, chunking, indexing, retrieval, reranking, context assembly, prompt construction, generation, citation, verification, and sometimes tool use or multi-turn interaction. Each stage can introduce reliability errors. A noisy document can enter during acquisition; chunking can split evidence away from its qualifiers; dense retrieval can surface semantically similar but answer-irrelevant passages; reranking can over-prioritise fluent misinformation; prompt construction can place misleading evidence before correct evidence; and generation can over-trust the retrieved context.

\paragraph{Evidence reliability dimensions.}
We use \emph{evidence reliability} as an umbrella concept with six dimensions. \textbf{Relevance} asks whether the evidence is about the query. \textbf{Sufficiency} asks whether it contains enough information to answer. \textbf{Faithfulness support} asks whether the final claim is entailed by the evidence rather than merely related. \textbf{Consistency} asks whether different evidence sources agree. \textbf{Provenance and freshness} ask where the evidence came from and whether it is up to date. \textbf{Integrity} asks whether the evidence has been manipulated. Existing work usually studies one or two of these dimensions at a time. For example, irrelevant-context robustness focuses on relevance and distraction \citep{irrelevantContext2023}, sufficient-context work focuses on answerability \citep{sufficientContext2024}, conflict work focuses on consistency \citep{tugOfWar2024,faithfulRAG2025}, and poisoning work focuses on integrity \citep{poisonedRAG2024,badRAG2024}. A central goal of this survey is to place these dimensions in one framework.

\paragraph{Benign unreliability vs. adversarial unreliability.}
Some evidence failures are benign: a retriever misses a relevant document, a query is ambiguous, or a corpus contains stale information. Others are adversarial: an attacker plants poisoned passages, modifies prompt templates, or injects hidden instructions into retrievable content. The two regimes share mechanisms. A hard negative and a poisoned passage can both be semantically close to the query; a stale document and a malicious document can both conflict with parametric knowledge; an ambiguous query and an adversarial trigger can both cause retrieval instability. However, adversarial settings require stronger assumptions about attacker capability, corpus mutability, and defense bypass. This is why we treat adversarial corruption as a distinct section while still connecting it to noisy retrieval and conflict.

\paragraph{Why robust RAG differs from robust LLMs.}
Robustness in closed-book LLMs often focuses on input perturbations, calibration, jailbreaks, or distribution shift. RAG adds an external memory whose content changes independently of the model. This creates three asymmetries. First, the corpus can be updated or poisoned without retraining the model. Second, the retriever can fail even when the generator is strong. Third, the generator may treat retrieved text simultaneously as evidence, instruction, and context. Therefore, robust RAG requires reasoning about the interaction between retriever, corpus, prompt, and generator, not only the generator.

\paragraph{Historical foundations and boundary.}
Many ingredients of robust RAG predate the current RAG literature. Retrieval-augmented pretraining, dense retrieval, and fusion-in-decoder open-domain QA established the retrieval--reader/generator interface \citep{realm2020,dpr2020,fid2021}. Selective prediction and neural calibration provide the conceptual basis for coverage--risk and confidence calibration \citep{geifman2017selective,guo2017calibration}. Long-context work shows that simply placing relevant evidence in a long input does not guarantee use of that evidence \citep{liu2024lostmiddle}. These works are background foundations rather than members of the venue-augmented robust-RAG evidence map unless they directly study RAG evidence reliability. We cite them to avoid presenting older IR, QA, calibration, and context-utilisation problems as entirely new.

\paragraph{A running example.}
Consider a hypothetical clinical-document assistant answering whether a drug is safe for a patient group. A benign failure might retrieve an outdated guideline or a passage about a similar but different drug. A distraction failure might retrieve a high-scoring passage that mentions the drug and patient group but discusses a different contraindication. A conflict failure might retrieve two guidelines with different dates. An adversarial failure might retrieve a poisoned webpage recommending unsafe use. An uncertainty failure might occur when no retrieved passage covers the subgroup. This example is illustrative rather than evidence that the surveyed benchmarks validate clinical deployment. It shows why robust RAG is a control problem grounded in evidence reliability: the system may need to re-rank, check dates, compare sources, abstain, or ask for more information rather than simply answer.
